
\chapter{Metodologia e Arquitetura do Sistema}
\label{chap:metodologia-arquitetura}

% Apresentar a abordagem metodológica adotada no trabalho.
% Explicar como o desenvolvimento do DoseCerta foi organizado
% e introduzir a arquitetura geral e a integração dos módulos.

\section{Caracterização da pesquisa}
\label{sec:caracterizacao-pesquisa}

% Caracterizar a natureza e os objetivos da pesquisa,
% bem como os procedimentos técnicos adotados.
% Justificar a classificação com base no trabalho realizado.

\section{Etapas de desenvolvimento}
\label{sec:etapas-desenvolvimento}

% Apresentar a sequência de etapas utilizadas no desenvolvimento.

\subsection{Levantamento de necessidades}
\label{subsec:levantamento-necessidades}

% Descrever como foram identificadas as necessidades
% do público-alvo e do sistema.

\subsection{Definição e especificação dos requisitos}
\label{subsec:especificacao-requisitos}

% Explicar como os requisitos funcionais e não funcionais
% foram definidos e documentados.

\subsection{Projeto e definição da arquitetura}
\label{subsec:projeto-arquitetura}

% Descrever como foram definidos os componentes do sistema,
% suas responsabilidades e as interfaces entre eles.

\section{Visão geral e fluxo de integração}
\label{sec:visao-geral-integracao}

% Apresentar a arquitetura geral do DoseCerta.
% Inserir um diagrama de blocos que represente os subsistemas
% e suas conexões, caso disponível.
%
% Explicar o funcionamento do sistema de ponta a ponta:
% interação do usuário, processamento das informações,
% comunicação entre módulos e operação do dispensador,
% conforme a implementação realizada.

\subsection{Módulos e responsabilidades}
\label{subsec:modulos-responsabilidades}

% Apresentar brevemente a função de cada módulo:
% projeto mecânico, eletrônica, sistema embarcado,
% backend, aplicativo, visão computacional e agente de IA,
% incluindo somente os módulos que fazem parte da solução.

\subsection{Fluxo de dados e comunicação}
\label{subsec:fluxo-dados-comunicacao}

% Explicar como os dados e comandos circulam entre os módulos.
% Informar interfaces e protocolos utilizados, quando aplicável.
% Remeter aos capítulos específicos para os detalhes técnicos.

\section{Implementação e integração dos componentes}
\label{sec:implementacao-integracao}

% Descrever a abordagem de implementação dos subsistemas
% e a sequência utilizada para integrá-los.
% Os detalhes técnicos de cada componente são apresentados
% nos capítulos específicos do projeto.

\subsection{Desenvolvimento dos subsistemas}
\label{subsec:desenvolvimento-subsistemas}

% Resumir como foram desenvolvidos os componentes mecânicos,
% eletrônicos, embarcados e de software, conforme aplicável.

\subsection{Integração do sistema}
\label{subsec:integracao-sistema}

% Explicar como os componentes foram integrados
% e quais procedimentos foram utilizados para verificar
% a comunicação e o funcionamento conjunto.

\section{Materiais, ferramentas e tecnologias}
\label{sec:materiais-ferramentas}

% Apresentar os materiais, equipamentos, plataformas,
% linguagens, bibliotecas e ferramentas utilizados.
% Organizar por categoria, se necessário.

\subsection{Hardware e componentes físicos}
\label{subsec:hardware}

% Listar os componentes físicos utilizados e suas funções.
% Os detalhes de projeto e especificações podem ser apresentados
% nos capítulos específicos.

\subsection{Software e ferramentas de desenvolvimento}
\label{subsec:software}

% Listar as ferramentas de software utilizadas
% e explicar sua aplicação no desenvolvimento.

\section{Procedimentos de teste e avaliação}
\label{sec:procedimentos-avaliacao}

% Descrever a estratégia de avaliação do sistema.
% Informar cenários de teste, critérios de avaliação
% e métodos de registro dos resultados.
% Os testes executados e seus resultados são detalhados
% no Capítulo 14.

\subsection{Testes funcionais}
\label{subsec:testes-funcionais}

% Descrever os procedimentos destinados a verificar
% se as funcionalidades atendem aos requisitos.

\subsection{Avaliação de usabilidade}
\label{subsec:avaliacao-usabilidade}

% Descrever os procedimentos de avaliação de usabilidade,
% caso tenham sido realizados ou planejados.
% Informar participantes, tarefas e critérios, quando aplicável.

\section{Limitações metodológicas}
\label{sec:limitacoes-metodologicas}

% Registrar limitações relacionadas aos recursos disponíveis,
% ao escopo, às condições de teste ou aos procedimentos
% de avaliação, quando aplicável.
